\section{Applications and negative-to-positive synthesis}

The experiments are deliberately small, but their failure modes suggest a practical design discipline. We distinguish three evidence levels throughout this section. A \emph{supported design implication} follows from the bounded P9 evidence itself; a \emph{plausible application} is motivated by primary-source systems that expose structurally analogous state; and a \emph{prospective pilot} is a new experiment that remains unrun. None of the application examples below is counted as a P9 performance result.

\subsection{Negative findings as design information}

P9 does not treat a negative result as an instruction to make the model larger. Instead, each negative is assigned a responsibility and converted into the minimum successor hypothesis that could repair it.

\paragraph{Topology without relation meaning.}
The D0 relation hostile pairs are non-identifying when only incidence/topology is visible. This negative yields the supported implication that a graph edge is not automatically the relevant scientific relation. When relation meaning is externally grounded, a system should expose that meaning explicitly rather than asking a deeper message-passing stack to infer information that is absent from its input. The positive successor hypothesis is therefore \emph{typed-relation sufficiency}, not generic graph depth.

\paragraph{Typed relations without transformation data.}
The transport pairs show that typed endpoints can still be non-identifying when the actual local transformation is hidden. The bounded A5 result then turns this limitation into a positive hybrid rule: once exact transformation values are visible, deterministic composition should receive first right of refusal. Learned models are better placed upstream, where a map must be proposed or estimated from noisy observations, than downstream, where a known map can be composed exactly.

\paragraph{Current state without relevant history.}
The failure-history pairs show that two current states may be model-identical while the productive action differs because an admitted prior failure is visible only in the semantic history view. The positive successor hypothesis is not ``remember every failure.'' It is that \emph{scoped, context-compatible negative history} can prevent repeated dead ends while stale or incompatible failures must not globally poison a mechanic. The corrected A2/A4 executable result remains pending and therefore this is still a bounded hypothesis outside the already-proved information ceiling.

\paragraph{No neural requirement on an exact atom.}
A5 is a negative result for neural escalation but a positive result for system design. A hybrid architecture may use learning for uncertain perception, extraction, ranking, or map proposal while reserving known invariants and closed-form constraints for explicit inference. On structure-native tasks, this separation may improve auditability and resource use while making missing information fail closed to \textsc{unknown}. A dedicated application pilot is required before claiming end-to-end benefit.

\paragraph{No trustworthy broad natural-paper gold.}
Closing the data lane at exact/procedural traces rather than fabricating natural-science labels is also a positive deployment constraint. P9 should first be tested in domains where state, transitions, dependencies, and validation are native artifacts rather than retrospective annotations. This motivates software incidents/program repair, formal proof, and process/workflow diagnosis before broad natural-science extraction.

\paragraph{Deferred latent, binding, and causal machinery.}
P9's bounded paper does not experimentally require recurrent latent compute, explicit role/entity binding, or intervention-stable causal mechanics. This is not evidence that those mechanisms are useless. The positive rule is \emph{triggered complexity}: reopen a mechanism only when a protected application contains a discriminator that simpler state, learning, or explicit inference cannot solve.

\subsection{Application 1: software incident diagnosis and program repair}

Software reliability is the highest-priority structure-native application because dependencies, traces, failures, candidate repairs, and external validation are already explicit artifacts. ITBench provides reproducible IT-automation scenarios and interpretable root-cause metrics; recent graph-guided root-cause work reasons over typed evidence graphs while deterministic graph/tool operations collect evidence and validate verdicts. Recent program-repair work similarly combines code-property and temporal-execution graphs with independent evidence perspectives rather than relying only on raw execution text.

A P9 pilot should therefore freeze incident families with same-symptom/different-root-cause controls, remint service/function identifiers, separate current evidence from admitted failed repairs, and include cases where decisive telemetry is absent so \textsc{unknown} is correct. The primary comparison should be topology-only, typed evidence graph, same-information serialization, typed graph plus scoped failure history, and typed graph plus explicit traversal/checking. End-to-end LLM remediation belongs to P10; P9 should first ask whether the structural state itself improves diagnosis and failure localization under matched information.

\subsection{Application 2: formal mathematics and proof planning}

Formal proof supplies exact states, explicit candidate operators, dependency structure, and a native verifier. Graph-based Lean systems such as Nazrin already demonstrate that expression graphs and atomic tactics are useful computational objects, so P9 does not claim graph theorem proving as novelty. Instead, formal proof is a strong test bed for the P9 decomposition: which proof obligations are non-identifying from a restricted view, when candidate tactics become identifiable after typed dependencies are exposed, when a failed tactic history changes the productive next action, and when an exact checker/search procedure should replace learned inference. Any claim that P9 improves theorem-proving success belongs to a separately frozen result, potentially P10 when an LLM/tool interface is involved.

\subsection{Application 3: process and workflow diagnosis}

Transactional workflows already appear as the protected D1 held-out domain. Object-centric process mining strengthens the real-world application case because modern event-log standards explicitly represent events, multiple object types, object--object and event--object relations, changing attributes, and qualified roles. Process-mining work also shows that flattening such processes can create artificial anomalies, while conformance-based methods can localize deviations to model elements. This is structurally close to P9's distinction between surface sequence, topology, typed relations, and explicit obstruction.

A prospective workflow pilot should use object-centric event logs with held-out process families, qualified relations, structural near-misses, explicit missing-information cases, and same-information serialized controls. The positive target is not generic anomaly detection; it is whether typed relational state plus exact conformance/inference improves localization and calibrated unresolved behavior over flattened or untyped views.

\subsection{Application 4: robotics and task planning as a triggered second wave}

Robotics is a natural later application because task-and-motion planning already represents symbolic operators through preconditions and effects, while learned skills introduce uncertain effects and role/entity grounding. Prior work on learning symbolic operators for task-and-motion planning shows that explicit operator abstractions can outperform model-free graph baselines on long-horizon planning, and newer neuro-symbolic systems combine learned skills with logical planning components. These donors remove any claim that P9 invented operator learning or hybrid planning.

Robotics is nevertheless valuable as a \emph{trigger test}. If a protected robotics task requires binding one reusable operator template to novel entities, A9 should be reopened. If observationally similar trajectories produce different effects under intervention, A10 should be reopened. If adaptive recurrent computation remains necessary after explicit state and search are matched, A3 should be reopened. This makes deferred complexity empirically demand-driven rather than architectural decoration.

\subsection{A real-world deployment hypothesis}

The strongest practical hypothesis suggested by P9 is therefore a hybrid rather than a monolithic network:
\[
\text{raw evidence}
\;\xrightarrow{\text{learned extraction/proposal}}\;
\text{typed epistemic state}
\;\xrightarrow{\text{learned ranking where needed}}\;
\text{exact/search inference where available}
\;\xrightarrow{\text{external validation}}\;
\text{result or unresolved}.
\]
The model is used where uncertainty and generalization require learning; exact computation is used where the state already supplies the required operands; and \textsc{unknown} is preserved when the decisive coordinate is absent. The open empirical question is whether this decomposition improves verified real-world outcomes, resource use, or failure localization. Issue~\#536 freezes the application research programme; no such application advantage is claimed by the present manuscript without a new result-bearing protocol.
